\chapter*{译者说明}
\markboth{译者说明}{}
\addcontentsline{toc}{chapter}{译者说明}

这份译注本译自 D. Gaitsgory 的预印本 \emph{Notes on Geometric Langlands: The Extended Whittaker Category}（2010 年 12 月 2 日，17 页）。它是作者「几何朗兰兹笔记」系列中的一册，该系列以 [GL:…] 为代号互相引用，本文多处结论的证明留在尚未写出的后续册中，作者用「Quasi-」（拟）标出。译文保留这一区分：凡原文标 Quasi- 者，译作「拟定理」「拟命题」等，不因中文读来生硬而抹平。

\medskip
\noindent\textbf{几点体例}

\begin{enumerate}[leftmargin=2.2em, itemsep=0.35em]
\item 原书的节号（1--8）、编号段（1.1、2.1.1……）与定理类编号（1.3.1、2.2.1……）共用一套数字，译文照旧，不另编号。
\item 原书的映射用 Fraktur 体 $\fP,\fQ,\fI,\fJ,\fR$，加横线者为紧化版本（$\bfP,\bfQ$），再加上标 $rat,red,ext,spec$ 层层区分；这些记号在中文排版中极易看错，译文一律原样保留，并在首次出现处加括注。
\item 原文的笔误（如 \texttt{consant}、$\mathrm{Av}_{\mathcal H_{G,H}}$、\texttt{neither zeros not poles}、定理编号误写、§3.2 两处恒等式的函子方向不合）均按「保留原文、另加译注」处理，不径改原文。各条的理由见工程文件 \texttt{work/疑点清单.tsv}。
\item 每条解释性译注都标注来源：\textbf{A} 指原书、作者他文或经典教材，\textbf{B} 指 nLab、Stacks Project、ProofWiki 一类公开辞书（须与 A 交叉），\textbf{I} 指本仓库已交付工程的术语表；凡无外部依据、仅由上下文推断者，标「译者按（上下文推断，无外部依据）」。
\item 术语依工程文件 \texttt{work/glossary.tsv}（冻结版）统一；「既约（reduced）」「系数函子」「Poincaré 级数函子」等取舍的理由见该表与 \texttt{work/通读笔记.md}。
\end{enumerate}

\medskip
\noindent\textbf{原书页与本书节的对照}

\begin{center}
\begin{tabular}{lll}
\toprule
原书页 & 内容 & 本书 \\
\midrule
p.1 & 引言 & 引言 \\
p.1--2 & §1 到 Borel 的有理约化空间 & §1 \\
p.2--4 & §2 施加 $N(\A)$-等变 & §2 \\
p.4--5 & §3 Eisenstein 级数函子 & §3 \\
p.6--8 & §4 谱图景 & §4 \\
p.8--10 & §5 Whittaker 范畴 & §5 \\
p.10--11 & §6 既约 Whittaker 范畴 & §6 \\
p.12--16 & §7 扩展 Whittaker 范畴 & §7 \\
p.16--17 & §8 粘合与谱图景 & §8 \\
p.17 & 参考文献 & 文献 \\
\bottomrule
\end{tabular}
\end{center}

\medskip
\noindent 底本为作者未正式发表的预印本，权利归原作者；本译本仅作内部研读之用，不对外发布。
